% Generated draft: requires mathematical and source review.
\documentclass[11pt]{article}
\usepackage[margin=25mm]{geometry}
\usepackage{fontspec}
\setmainfont{Latin Modern Roman}
\usepackage{amsmath,amssymb,mathtools,unicode-math}
\setmathfont{Latin Modern Math}
\usepackage{enumitem,needspace,xurl,hyperref}
\setlength{\parindent}{0pt}
\setlength{\parskip}{5pt}
\setlength{\emergencystretch}{3em}
\newcommand{\N}{\mathbb{N}}
\newcommand{\Z}{\mathbb{Z}}
\newcommand{\Q}{\mathbb{Q}}
\newcommand{\R}{\mathbb{R}}
\newcommand{\C}{\mathbb{C}}
\newcommand{\F}{\mathbb{F}}
\newcommand{\A}{\mathbb{A}}
\newcommand{\PP}{\mathbb{P}}
\newcommand{\E}{\mathbb{E}}
\newcommand{\eps}{\varepsilon}
\newcommand{\dd}{\,\mathrm{d}}
\begin{document}
\section*{UnsolvedMath ID: 800311}
\textbf{Source problem code:} catalogue entry
\subsection{Definitions and mathematical statement}
The supplied question is: ``Is the (restriction of the) Godbillon--Vey class invariant under topological conjugacy for groups of \(C^2\) diffeomorphisms?''

Under a group-cohomological interpretation of ``restriction'', its invariance assertion has the following form. Let \(M\) be the underlying one-dimensional manifold, let \(D\) be the intended group of \(C^2\) diffeomorphisms of \(M\), and let
\[
\operatorname{GV}_M\in H^q(D^\delta;\R)
\]
be the intended Godbillon--Vey class. Here \(D^\delta\) denotes the group equipped with the discrete topology, and \(H^q\) denotes group cohomology with trivial real coefficients. The manifold, the permitted diffeomorphisms, and the precise class and degree must be specified before this schema determines a particular conjecture.

For subgroups \(G_1,G_2\leq D\), write \(\iota_j:G_j\hookrightarrow D\) for the inclusions and define the restricted classes by
\[
\operatorname{GV}_{G_j}=\iota_j^*\operatorname{GV}_M
\in H^q(G_j^\delta;\R).
\]
A topological conjugacy is a homeomorphism \(h:M\to M\) such that \(hG_1h^{-1}=G_2\); it induces the group isomorphism
\[
c_h:G_1\longrightarrow G_2,\qquad c_h(g)=h\circ g\circ h^{-1}.
\]
The question is whether, for every such pair of subgroups and every admissible topological conjugacy, one has
\[
 c_h^*\operatorname{GV}_{G_2}=\operatorname{GV}_{G_1}
 \quad\text{in }H^q(G_1^\delta;\R).
\]
The conjugating homeomorphism is not assumed to be differentiable. The supplied record does not specify the underlying manifold, boundary or support conditions, orientation conventions for the diffeomorphisms and conjugacy, or the cohomological realization of the class. Consequently, the displayed assertion is an explicit conditional formulation of the restriction-invariance question, rather than a uniquely specified conjecture extracted from the record. In particular, whether orientation-reversing conjugacies are included, and whether they require a sign convention, remains unresolved.
\subsection{Short English statement}
Does topologically conjugating a group of \(C^2\) diffeomorphisms preserve its restricted Godbillon--Vey class, after identifying the groups by conjugation? The supplied record leaves the manifold, orientation conditions, and precise realization of the class unspecified. Its supplied status is \texttt{open}; this draft does not independently verify that assessment.
\subsection{Sources}
\begin{itemize}
\item Local UnsolvedMath registry, supplied complete source record: permanent ID \(800311\), problem code \texttt{catalogue entry}, title ``Group Actions on 1-Manifolds --- Question 28''. Attribution: Navas, proposed in 2017. The record supplies the status \texttt{open} and an update timestamp of \texttt{2026-10-16T14:00:00Z}.
\item Navas, \emph{Group Actions on 1-Manifolds: A List of Very Concrete Open Questions} (2017), supplied source locator ``Question 28 (\#27)'': \url{https://arxiv.org/abs/1712.06462}. The supplied background records access on \texttt{2026-10-16} and extraction by \texttt{arxiv\_latex}. This access information is supplied provenance, not an independently verified visit or literature assessment.
\end{itemize}
\end{document}
